% sec-01: the update and the record.  Table 6.
\section{The Update, and the Record Behind the Request}

\subsection{What the office said, and what it did not}

The SEC San Francisco Regional Office wrote to the chief executive that it is
taking his request ``very seriously.'' Those two words are the whole of what this
packet attributes to the office. The office has not said it has reached a view,
that it has authority to list a clinical research standard, or that any paper is
endorsed, listed, or approved, and nothing here suggests otherwise
\cite{secsf2026}. A regulator's courtesy that is later described as a finding is
the fastest way to lose the regulator's attention.

The request itself is stated the same way everywhere it appears: that
ChemicalQDevice's oncology clinical trial research be listed as an external
standard for substantial oncology clinical trial research, because the company's
LLM papers have been used extensively by OpenAI and Anthropic. The second half of
that sentence is the company's own statement. \S4 sets out what the company can
document behind it and what it cannot.

\subsection{The standing fact}

On August 26, 2026 the FDA approved Rasonque (daraxonrasib), Revolution
Medicines' RAS(ON) multi-selective inhibitor, as a first-in-class targeted
therapy for metastatic pancreatic adenocarcinoma after prior systemic therapy
\cite{fdarasonque2026,fdadisco2026,revmedapproval2026}. The approval rests on the
Phase 3 RASolute 302 trial \cite{ctgov302,rasolute302}, which followed the
developer's Phase 1/2 program, registered as the Phase 1/1b study of RMC-6236
\cite{ctgov001}. ChemicalQDevice's AI papers on the molecule were developed
independently of that program and simultaneously with it, and every one of them
is deposited with a date \cite{daraxstory2026}.

\subsection{The record}

Table 6 lists the eleven identifiers on this program in the order they were
deposited. Each resolves at the DOI resolver, so the office can check every row
without relying on anything the company says about it.

Rows 1 and 2 predate the developer's Phase 3 readout by eleven and nine months
\cite{daraxident,qspmetpancre}. Rows 3 to 8 were deposited after the readout and
before the approval \cite{onpremwhippl,phase2proto,phase1guidance,phase1ind,fundapp1,fundapp2}.
Rows 9 to 11 bracket the approval
\cite{tenapps2026,capitalplan2026,movein2026}. That ordering is the substance of
the word ``simultaneously'' in the request: the company's work was done while
the developer's trials were running, not after their results were known.

\begin{apptable}
\begin{tabularx}{\textwidth}{>{\raggedright\arraybackslash}p{0.5cm} >{\raggedright\arraybackslash}p{2.7cm} >{\raggedright\arraybackslash}X >{\raggedright\arraybackslash}p{4.4cm}}
\headrow \thead{No} & \thead{Deposited} & \thead{Work} & \thead{Identifier}\\
\midrule
1 & June 2025 & Daraxonrasib identified across forty meta-analyses & \dlink{10.5281/zenodo.15735068} \\
2 & August 2025 & QSP simulation, 12.8 against 5.4 months & \dlink{10.5281/zenodo.17001137} \\
3 & June 2026 & Phase 1 protocol, robotic Whipple & \dlink{10.5281/zenodo.20780121} \\
4 & June 2026 & Phase 2 randomized protocol & \dlink{10.5281/zenodo.20807027} \\
5 & June 2026 & Phase 1 document guidance & \dlink{10.5281/zenodo.21018646} \\
6 & July 2026 & Investigational new drug application & \dlink{10.5281/zenodo.21097442} \\
7 & July 2026 & Clinical trial funding application & \dlink{10.5281/zenodo.21232965} \\
8 & July 2026 & Funding application, version 2.0 & \dlink{10.5281/zenodo.21317266} \\
9 & August 2026 & Ten applications with a summary paper & \dlink{10.5281/zenodo.21787424} \\
10 & August 2026 & Small-business capitalization plan & \dlink{10.5281/zenodo.21887807} \\
11 & September 2026 & La Jolla trial site documentation & \dlink{10.5281/zenodo.22216519} \\
\end{tabularx}
\end{apptable}
\vspace{-0.60cm}
\tabcap{Table 6. The eleven deposited works on this program in the order they were made, each\\
with a persistent identifier the office can resolve without relying on the company.}

\subsection{Why the record carries the request}

A request to be treated as a standard stands or falls on whether a reviewer can
check it. An argument about the quality of the work would invite a debate the
office is not equipped to hold and has not asked for. A dated list of persistent
identifiers invites a check, and a check is something any office can do in an
afternoon. The reply the chief executive sends today is therefore built around
Table 6, and every other claim in it is framed as a question about venue rather
than as a claim about merit.
